\epanalipsi{1}

\tmima{Λεξικό της ενότητας}

\renewcommand{\arraystretch}{1.5}
\begin{tabular*}{\linewidth}{@{\extracolsep{\fill}} >{\raggedright\arraybackslash}p{0.26\linewidth} >{\raggedright\arraybackslash}p{0.64\linewidth}}
\toprule
\textbf{Λέξη ή τύπος} & \textbf{Με μία πρόταση}\\
\midrule
φυσικός & Ένας από τους $0,\,1,2,\,3,\dots$: μετράει πράγματα που δεν σπάνε.\\
στρογγυλοποίηση & Αλλάζω έναν αριθμό με κοντινό του, με συμφωνημένα βήματα, μία φορά.\\
$\alpha^{\nu}$ & Ο $\alpha$ γραμμένος $\nu$ φορές σε γινόμενο.\\
προτεραιότητα & Παρενθέσεις, δυνάμεις, $\cdot$ και $:$, και τέλος $+$ και $-$.\\
$\Delta=\delta\cdot\pi+\upsilon$ & Ευκλείδεια διαίρεση, με $\upsilon<\delta$.\\
κριτήρια & Κανόνες στα ψηφία που δείχνουν αν ένας αριθμός διαιρείται με το 2, 3, 4, 5, 9, 10, 25.\\
πρώτος & Ο αριθμός μεγαλύτερος από το 1 με διαιρέτες μόνο τον 1 και τον εαυτό του.\\
ΜΚΔ & Ο μεγαλύτερος κοινός διαιρέτης: μοιράζω.\\
ΕΚΠ & Το μικρότερο (μη μηδενικό) κοινό πολλαπλάσιο: συναντιέμαι.\\
\bottomrule
\end{tabular*}

\tmima{Ασκήσεις επανάληψης}

\ex Γράψε με ψηφία τον αριθμό «τριάντα επτά εκατομμύρια, δύο χιλιάδες, πέντε». Μετά στρογγυλοποίησέ τον στα εκατομμύρια.

\ex Βάλε σε αύξουσα σειρά τους 5.012, 5.102, 5.021, 4.999 και 10.000.

\ex Στρογγυλοποίησε το 86.549 στις δεκάδες, στις εκατοντάδες και στις χιλιάδες.

\ex Υπολόγισε με τις ιδιότητες των πράξεων: α) $25\cdot 17\cdot 4$ \quad β) $99\cdot 8$ \quad γ) $36\cdot 15+36\cdot 85$ \quad δ) $47+68+53+32$.

\ex Βρες την τιμή: α) $3\cdot(8-2)^2-10$ \quad β) $2^5-5^2$ \quad γ) $(2+3)\cdot 4-2^3$ \quad δ) $100-4\cdot 5^2$.

\ex Κάνε την Ευκλείδεια διαίρεση $1.000:48$ και γράψε την ισότητα $\Delta=\delta\cdot\pi+\upsilon$. Ποιος αριθμός δίνει πηλίκο 7 και υπόλοιπο 6 όταν διαιρεθεί με το 15;

\ex Σήμερα είναι Πέμπτη. Τι μέρα θα είναι μετά από 100 ημέρες;

\ex Με ποιους από τους 2, 3, 4, 5, 9, 10 και 25 διαιρείται ο αριθμός 3.240;

\ex Ποιο ψηφίο πρέπει να μπει στη θέση του $\square$ ώστε ο $5\square 7$ να διαιρείται με το 9;

\ex Ποιοι από τους 57, 59, 63 και 67 είναι πρώτοι; Για τους σύνθετους, βρες την ανάλυσή τους.

\ex Ανάλυσε τον 360 σε γινόμενο πρώτων παραγόντων. Πόσους πρώτους διαφορετικούς παράγοντες έχει;

\ex Βρες το ΜΚΔ και το ΕΚΠ των αριθμών 36 και 60.

\ex \textbf{Οι κάρτες του Στέλιου.} Ο Στέλιος έχει 36 κάρτες με ζώα και 60 κάρτες με αυτοκίνητα. Θέλει να φτιάξει όσο το δυνατόν περισσότερες ίδιες στοίβες, με και από τα δύο είδη, χωρίς να περισσέψει καμία. Πόσες στοίβες θα γίνουν και τι έχει η κάθε μία;

\ex \textbf{Η πισίνα.} Ο Ηλίας πάει στην πισίνα κάθε 8 μέρες και η Ζωή κάθε 12. Σήμερα πήγαν μαζί. Πότε θα ξαναπάνε μαζί;

\ex Σωστό ή λάθος; α) Κάθε άρτιος αριθμός διαιρείται με το 4. \quad β) $\MKD(12,\,18)=6$. \quad γ) Το 91 είναι πρώτος. \quad δ) $2^4=8$. \quad ε) $17-(9-3)=17-9-3$. \quad στ) $0:5=0$.

\ex Ο Άρης λέει ότι $6+4\cdot 5=50$ και η Ζωή ότι $6+4\cdot 5=26$. Ποιος έχει δίκιο και γιατί; Πού πρέπει να βάλεις παρένθεση ώστε να βγει 50;

\ex Πόσους διαιρέτες έχει το 100; Πόσοι από αυτούς είναι πρώτοι;

\tmima{Ο χάρτης της επανάληψης}

Μετά τις λύσεις, δες σε ποιες ασκήσεις δυσκολεύτηκες. Ο πίνακας σου λέει ποιο κεφάλαιο να ξαναδιαβάσεις.

\renewcommand{\arraystretch}{1.5}
\begin{tabular*}{\linewidth}{@{\extracolsep{\fill}} >{\raggedright\arraybackslash}p{0.38\linewidth} >{\raggedright\arraybackslash}p{0.52\linewidth}}
\toprule
\textbf{Αν δυσκολεύτηκες στις} & \textbf{ξαναδιάβασε τα κεφάλαια}\\
\midrule
Ε1.1, Ε1.2, Ε1.3 & 1 (φυσικοί, διάταξη, στρογγυλοποίηση)\\
Ε1.4 & 2 (ιδιότητες των πράξεων)\\
Ε1.5, Ε1.16 & 3 (δυνάμεις, προτεραιότητα)\\
Ε1.6, Ε1.7 & 4 (Ευκλείδεια διαίρεση)\\
Ε1.8, Ε1.9 & 5 (κριτήρια διαιρετότητας)\\
Ε1.10, Ε1.11, Ε1.17 & 6 (πρώτοι, ανάλυση)\\
Ε1.12, Ε1.13, Ε1.14 & 7 (ΜΚΔ, ΕΚΠ)\\
Ε1.15 & 1 έως 7\\
\bottomrule
\end{tabular*}
